\section{Introduction}
\subsection*{General}
\begin{frame}[label=contents_prog1]
  \frametitle{Today's outline}
    \only<1>{\tableofcontents[hideallsubsections]}
    \only<2>{\tableofcontents[hideallsubsections,currentsection]}
\end{frame}

\section{Python tutorial slides}
\begin{frame}[fragile]
   \frametitle{Lists in Python (1)}
  \begin{itemize}
    \item Lists are containers of collections of objects
    \item A list is initialized using square brackets with comma-separated elements
    \begin{lstlisting}[language=Python,numbers=none]
>>> brands = ['Audi', 'Toyota', 'Honda', 'Ford', 'Tesla']
    \end{lstlisting}
    \item Lists can contain and mix any object type, even other lists:
    \begin{lstlisting}[language=Python,numbers=none]
>>> another_list = [0.0, 0.1, 0.2, 'Hello', brands]
>>> print(another_list)
    \end{lstlisting}
    \begin{lstlisting}[style=PyOutput]
[0.0, 0.1, 0.2, 'Hello', ['Audi', 'Toyota', 'Honda', 'Ford', 'Tesla']]
    \end{lstlisting}
    \item Access (i.e., read) an entry in a list. Note that indexing starts at 0:
    \begin{lstlisting}[language=Python,numbers=none]
>>> print(another_list[0],another_list[3])
    \end{lstlisting}
    \begin{lstlisting}[style=PyOutput]
0.0 Hello
    \end{lstlisting}
  \end{itemize}
 \end{frame}
 
 \begin{frame}[fragile]
   \frametitle{Lists in Python (2)}
   \begin{itemize}
    \item Manipulate the value of an entry goes likewise:
    \begin{lstlisting}[language=Python,numbers=none]
>>> another_list[3] = 'Bye' # Becomes: [0.0, 0.1, 0.2, 'Bye', ['Audi', ...]]
    \end{lstlisting}
    \item Slicing is used to retrieve multiple elements:
    \begin{lstlisting}[language=Python,numbers=none]
>>> another_list[1:4] # This will give the elements from index 1 to index 3
    \end{lstlisting}
    \begin{lstlisting}[style=PyOutput]
[0.1, 0.2, 'Bye']
    \end{lstlisting}
    \item Lists can be unpacked into individual variables:
    \begin{lstlisting}[language=Python,numbers=none]
>>> a,b,c,d,e = brands
>>> print(f"The first list element was {a}, then {b}, {c}, {d} and finally {e}.")
    \end{lstlisting}
    \begin{lstlisting}[style=PyOutput]
The first list element was Audi, then Toyota, Honda, Ford and finally Tesla.
    \end{lstlisting}
    \item From here onwards, we will omit the \lstinline|print| statements from the slides
   \end{itemize}\vskip1em
  \end{frame}
 
  \begin{frame}[fragile]
    \frametitle{Lists in Python (3)}
    \begin{itemize}
      \item Lists can be concatenated or repeated by the addition and multiplication operators respectively: 
      \begin{lstlisting}[language=Python,numbers=none]
>>> more_brands = ['Nissan','Kia'] + brands
      \end{lstlisting}
      \begin{lstlisting}[style=PyOutput]
['Nissan', 'Kia', 'Audi', 'Toyota', 'Honda', 'Ford', 'Tesla']
      \end{lstlisting}
      \begin{lstlisting}[language=Python,numbers=none]
>>> zeros = 10*[0]
      \end{lstlisting}
      \begin{lstlisting}[style=PyOutput]
[0, 0, 0, 0, 0, 0, 0, 0, 0, 0]
      \end{lstlisting}
      \item Find out which methods can be performed on a list by using \lstinline|dir(more_brands)|:
      \begin{lstlisting}
more_brands.append('Volvo') # Append object (here: string literal) at the end of the list 
more_brands.insert(1,'BMW') # Insert object at index 1
more_brands.sort()          # Sorts the list in-place
item = more_brands.pop(3)   # Removes element at index 3 from the list, stores it as item
      \end{lstlisting}
    \end{itemize}
\end{frame}
 
 \begin{frame}[fragile]
   \frametitle{Lists in Python (4)}
   Ranges of numbers are set using the \lstinline|range(start=0,stop,step=1)| command:
   \begin{itemize}
    \item Create a list with a range of numbers:
    \begin{lstlisting}[language=Python,numbers=none]
>>> a = list(range(1, 11))     # Creates a list from 1 to 10
    \end{lstlisting}
    \item List comprehensions can be used to create lists with more complex patterns:
    \begin{lstlisting}[language=Python,numbers=none]
>>> x = [i/10 for i in range(-10, 11)]   # Creates a list from -1 to 1 with a step of 0.1
    \end{lstlisting}
    \item Manipulating multiple components using slicing and a loop:
    \begin{lstlisting}[language=Python,numbers=none]
>>> y = list(range(11))  # Creates a list from 0 to 10
>>> for i in [0, 3, 4, 5, 6]: 
>>>     y[i] = 1
    \end{lstlisting}
    \item Or (by supplying a list instead of a scalar):
    \begin{lstlisting}[language=Python,numbers=none]
>>> y[0:2] = [16, 19]  # Sets y[0] to 16 and y[1] to 19
    \end{lstlisting}
  \end{itemize}
 \end{frame}

\begin{frame}[fragile]
  \frametitle{Assignment of variables; value or reference}
  Consider the following code snippets:
  \begin{columns}[t]
    \column{0.4\textwidth}
    \begin{lstlisting}
>>> i = 3
>>> j = i 
>>> j = i + 3
>>> print(i,j)
>>> print(id(i), id(j)) # Print memory address of data
    \end{lstlisting}\pause
    \begin{lstlisting}[style=PyOutput]
3, 6
8885416 8885544
    \end{lstlisting}\pause
    \column{0.5\textwidth}
    \begin{lstlisting}
>>> list_a = ['aa', 1, 'bb', 12, True, 1.618]
>>> list_b = list_a
>>> list_b[2] = 'cc'
>>> print(list_a, list_b, sep='\n')
>>> print(id(list_a), id(list_b))
    \end{lstlisting}\pause
    \begin{lstlisting}[style=PyOutput]
['aa', 1, 'cc', 12, True, 1.618]
['aa', 1, 'cc', 12, True, 1.618]
140285003056512 140285003056512
    \end{lstlisting}\pause
  \end{columns}
  \begin{itemize}
    \item Primitive or immutable data types (e.g. \lstinline|int|, \lstinline|float|, \lstinline|str|, \lstinline|tuple|) are \emph{assigned by value}; the value is copied and changes do not affect the original variables.
    \item Mutable data types (e.g. \lstinline|list|, \lstinline|set|, \lstinline|dict|) are \emph{assigned by reference}; they are two names pointing to the same data, changing one affects the values of the other.
  \end{itemize}
\end{frame}



%% STRING SECTION
\subsection{Strings}
\begin{frame}[fragile]
  \frametitle{Strings in Python (1)}
  Creating a string:
  \begin{lstlisting}[language=Python,numbers=none]
>>> s = "Hello, world!"
>>> len(s)
13
  \end{lstlisting}
  Accessing a character in a string:
  \begin{lstlisting}[language=Python,numbers=none]
>>> s[7]
'w'
  \end{lstlisting}
  Getting a substring:
  \begin{lstlisting}[language=Python,numbers=none]
>>> s[7:12]
'world'
  \end{lstlisting}
  Or separate by whitespace using a string method (see \lstinline|dir(s)|):
  \begin{lstlisting}[language=Python,numbers=none]
>>> s.split()
['Hello,', 'world!']
  \end{lstlisting}
\end{frame}

\begin{frame}[fragile]
  \frametitle{Strings in Python (2)}
  Replacing a substring with another string:
  \begin{lstlisting}[language=Python,numbers=none]
>>> s.replace('world', 'Python')
'Hello, Python!'
  \end{lstlisting}
  Converting to upper and lower case:
  \begin{lstlisting}[language=Python,numbers=none]
>>> s.upper()
'HELLO, WORLD!'
>>> s.lower()
'hello, world!'
  \end{lstlisting}
  You can combine methods with string literals too:
  \begin{lstlisting}[language=Python,numbers=none]
>>> s.replace('WoRlD'.lower(), 'Python')
'Hello, Python!'
>>> s.startswith('hello'.title())
True
  \end{lstlisting}
  Finding the starting index of a substring:
  \begin{lstlisting}[language=Python,numbers=none]
>>> s.index("world")
7
  \end{lstlisting}
\end{frame}

\begin{frame}[fragile,label=practice_strings]
  \frametitle{Practice}
  Given a string
  \begin{lstlisting}
>>> s = "Python programming is fun!"
  \end{lstlisting}
  \begin{itemize}
   \item Find and print the index of the word "is".
   \item Create a new string where "fun" is replaced with "awesome".
   \item Print the string in uppercase.
  \end{itemize}
  \begin{onlyenv}<beamer>
    \begin{lstlisting}[]
>>> s.index('is')
>>> p = s.replace('fun','awesome')
>>> print(p.upper())
    \end{lstlisting}
  \end{onlyenv}
 \end{frame}

 %% TUPLES SECTION
\subsection{Tuples}
\begin{frame}[fragile]
  \frametitle{Tuples in Python}
  A tuple is a built-in data type that contains an immutable sequence of values. 
  Creating a tuple:
  \begin{lstlisting}[language=Python,numbers=none]
>>> t = (1, 2, 3)
  \end{lstlisting}
  Accessing an element of a tuple:
  \begin{lstlisting}[language=Python,numbers=none]
>>> t[1]
2
  \end{lstlisting}
  Tuples are immutable, so we can't change their elements. However, we can create a new tuple based on the old one:
  \begin{lstlisting}[language=Python,numbers=none]
>>> t = t + (4, )
  \end{lstlisting}
  Finding the length of a tuple:
  \begin{lstlisting}[language=Python,numbers=none]
>>> len(t)
4
  \end{lstlisting}
\end{frame}

\begin{frame}[fragile,label=practice_tuples]
  \frametitle{Practice}
  Given a tuple
  \begin{lstlisting}
>>> t = (1, 2, 3, 4, 5, 6)
  \end{lstlisting}
  \begin{itemize}
   \item Access and print the third element of the tuple.
   \item Try to change the value of the second element of the tuple.
   \item Create a new tuple by concatenating a second tuple \((7, 8, 9)\) to the original tuple.
  \end{itemize}\pause
  \begin{onlyenv}<beamer>
    \begin{lstlisting}[]
t = (1, 2, 3, 4, 5, 6)
print(t[2])
t[2] = 6
t2 = t + (7,8,9)
print(t2)
    \end{lstlisting}
  \end{onlyenv}
 \end{frame}

%% DICTIONARY SECTION 
\subsection{Dictionaries}
\begin{frame}[fragile]
  \frametitle{Dictionaries in Python (1)}
  Creating a dictionary:
  \begin{lstlisting}[language=Python,numbers=none]
>>> d = {'a': 1, 'b': 2, 'c': 3}
  \end{lstlisting}
  Accessing a value by its key:
  \begin{lstlisting}[language=Python,numbers=none]
>>> d['b']
2
  \end{lstlisting}
  Modifying a value associated with a key:
  \begin{lstlisting}[language=Python,numbers=none]
>>> d['b'] = 47
  \end{lstlisting}
  Adding a new key-value pair:
  \begin{lstlisting}[language=Python,numbers=none]
>>> d['d'] = 4
  \end{lstlisting}
  Removing a key-value pair using pop:
  \begin{lstlisting}[language=Python,numbers=none]
>>> d.pop('d')
4
  \end{lstlisting}
\end{frame}

\begin{frame}[fragile]
  \frametitle{Dictionaries in Python (2)}
  Get all keys as a list:
  \begin{lstlisting}[language=Python,numbers=none]
>>> list(d.keys())
['a', 'b', 'c']
  \end{lstlisting}
  Get all values as a list:
  \begin{lstlisting}[language=Python,numbers=none]
>>> list(d.values())
[1, 47, 3]
  \end{lstlisting}
  Get all key-value pairs as a list of tuples:
  \begin{lstlisting}[language=Python,numbers=none]
>>> list(d.items())
[('a', 1), ('b', 47), ('c', 3)]
  \end{lstlisting}
\end{frame}

\begin{frame}[fragile,label=practice_dictionary]
  \frametitle{Practice}
  Given a dictionary
  \begin{lstlisting}[language=Python,numbers=none]
>>> d = { 'Alice': 24, 'Bob': 27, 'Charlie': 22, 'Dave': 30}
  \end{lstlisting}
  \begin{itemize}
   \item Access and print the age of 'Charlie'.
   \item Update 'Alice' age to 25.
   \item Add a new entry for 'Eve' with age 29.
   \item Print all the keys in the dictionary.
  \end{itemize}\pause
  \begin{onlyenv}<beamer>
    \begin{lstlisting}[]
>>> print(d['Charlie'])
>>> d['Alice'] = 25
>>> d['Eve'] = 29
>>> print(d.keys())
dict_keys(['Alice', 'Bob', 'Charlie', 'Dave', 'Eve'])
    \end{lstlisting}
  \end{onlyenv}
 \end{frame}

 \section{Control flow}
 \subsection{Loops}
 \againframe<2>{contents_prog1}
 
 %% Loops
\begin{frame}[fragile]
  \frametitle{Loops in Python (1)}
  The \lstinline$for$ loop is used to iterate over a sequence (e.g. lists, sets, tuples, dictionaries, strings). Any \emph{iterable} object can be listed over:
  \begin{lstlisting}[language=Python,numbers=none]
>>> for i in range(5):
...     print(i)
0
1
2
3
4
  \end{lstlisting}\pause
  You can iterate over a list directly:
  \begin{lstlisting}[language=Python,numbers=none]
>>> my_list = [1, 2, 3, 4, 5]
>>> for num in my_list:
...     print(num)
1
2
3
4
5
  \end{lstlisting}
\end{frame}

\begin{frame}[fragile]
  \frametitle{Loops in Python (2)}
  The \lstinline|enumerate| keyword returns both the \emph{index} as well as the \emph{list element}:
  \begin{lstlisting}[language=Python,numbers=none]
>>> my_list = ['aa', 1, 'bb', 12, True, 1.618034, []]
>>> for idx,elm in enumerate(my_list):
...     print(f'Element {elm} of type {type(elm)} at index {idx}')
  \end{lstlisting}
  \pause
  \begin{lstlisting}[style=PyOutput]
Element aa of type <class 'str'> at index 0
Element 1 of type <class 'int'> at index 1
Element bb of type <class 'str'> at index 2
Element 12 of type <class 'int'> at index 3
Element True of type <class 'bool'> at index 4
Element 1.618034 of type <class 'float'> at index 5
Element [] of type <class 'list'> at index 6
  \end{lstlisting}
\end{frame}

\begin{frame}[fragile]
  \frametitle{Loops in Python (3)}
  The `while` loop keeps going as long as a condition is \lstinline{True}:
  \begin{lstlisting}[language=Python,numbers=none]
>>> i = 0
>>> while i < 3:
...     print(i)
...     i += 1
0
1
2
  \end{lstlisting}\pause
  Use \lstinline|break| to exit a loop prematurely, and \lstinline|continue| to skip to the next iteration:
  \begin{lstlisting}[language=Python,numbers=none]
>>> for i in range(5):
...     if i == 3:
...         break
...     print(i)
0
1
2
  \end{lstlisting}
\end{frame}

\begin{frame}[fragile,label=practice_for]
  \frametitle{Practice}
  Given a list
  \begin{lstlisting}[language=Python,numbers=none]
>>> my_list = [1, 3, 7, 8, 9] 
  \end{lstlisting}
 
 \begin{itemize}
  \item Use a for loop to print each element of the list.\pause
  \item Use a while loop to print the values at indices 0 to 4.\pause
  \item Use a loop to find and print the index of the number 7 in the list.
 \end{itemize}
 \begin{onlyenv}<beamer>
  \begin{columns}[T]
    \column{0.33\textwidth}
      \begin{lstlisting}[]
>>> for i in my_list:
...    print(i)
      \end{lstlisting}
    \column{0.33\textwidth}
      \begin{lstlisting}[]
>>> c = 0
>>> while (c<4):
...    print(f"{my_list[c]=}")
...    c += 1
      \end{lstlisting}
    \column{0.33\textwidth}
      \begin{lstlisting}[]
>>> c = 0
>>> while not my_list[c] == 7:
...    c += 1
>>> print(my_list[c])
      \end{lstlisting}      
  \end{columns}
 \end{onlyenv}
\end{frame}

%% CONTROL STATEMENTS
\subsection{Branching}
\begin{frame}[fragile]
  \frametitle{Conditional Statements in Python}
  The Boolean type \lstinline|bool| has only 2 possible values: \lstinline|True| or \lstinline|False| \\
  The \lstinline{if} statement is used to execute a block of code only if a condition is evaluated to \lstinline|True|:
  \begin{lstlisting}[language=Python, numbers=none, deletekeywords={is,not}]
>>> x = 5
>>> if x > 0:
...     print("x is positive")
x is positive
  \end{lstlisting}\pause
  Use \lstinline{elif} to specify additional conditions, and \lstinline{else} to define what to do if no conditions are met:
  \begin{lstlisting}[language=Python, numbers=none, deletekeywords={is,not}]
>>> if x > 10:
...     print("x is greater than 10")
... elif x == 10:
...     print("x is exactly 10")
... else:
...     print("x is less than 10")
x is less than 10
  \end{lstlisting}
\end{frame}

\begin{frame}[fragile]
  \frametitle{Nested conditionals}
  Nesting conditions allows for more complex conditionals:
  \begin{lstlisting}[language=Python, numbers=none, deletekeywords={is,not}]
>>> if x > 0:
...     if x % 2 == 0: # The modulo operator % yields the remainder of a division
...         print("x is positive and even")
...     else:
...         print("x is positive but odd")
... else:
...     print("x is non-positive")
x is positive but odd
  \end{lstlisting}\pause
  The \lstinline{in} keyword can be used to check membership in a sequence:
  \begin{lstlisting}[language=Python, numbers=none, deletekeywords={is,not}]
>>> my_list = [1, 2, 3, 4, 5]
>>> if 3 in my_list:
...     print("3 is a member of the list")
3 is a member of the list
  \end{lstlisting}
\end{frame}

\begin{frame}[fragile]
  \frametitle{Combined conditionals}
  \begin{block}{Leap year determination}
    To be a leap year, the year number must be divisible by four, except for end-of-century years, which must be divisible by 400.
  \end{block}\pause
  We can create combined conditions using \lstinline|not|, \lstinline|and| and \lstinline|or| to determine whether we have a leap year:
  \pause
  \begin{lstlisting}[language=Python, numbers=none]
>>> year = 2024
>>> if year % 4 == 0 and (not year % 100 == 0 or year % 400 == 0):
...    print(f"{year} is a leap year.")
... else:
...    print(f"{year} is not a leap year.")
  \end{lstlisting}
\end{frame}

\begin{frame}[fragile]
  \frametitle{Conditionals in list comprehensions}
  List comprehensions can be extended with conditionals too:
  \begin{lstlisting}[language=Python, numbers=none]
>>> x = [i for i in range(0,31) if i%3 == 0]
>>> print(x)
[0, 3, 6, 9, 12, 15, 18, 21, 24, 27, 30]
  \end{lstlisting}\pause
  \vskip1em
  Conditions are not restricted to modulo's. Here we select artists who have an 'r' or 'R' in their name:
  \begin{lstlisting}[language=Python, numbers=none]
artists = ["Adele", "Harry Styles", "Stef Ekkel", "Ed Sheeran", "Nicki Minaj", "Ariana Grande", "Robbie Williams"]

my_artists = [artist for artist in artists if artist.lower().count('r') > 0]
print(my_artists)
['Harry Styles', 'Ed Sheeran', 'Ariana Grande', 'Robbie Williams']
  \end{lstlisting}
  
\end{frame}

\begin{frame}[fragile,label=practice_conditions]
  \frametitle{Practice}
  Consider the list:
  \begin{lstlisting}[language=Python, numbers=none]
>>> my_list = [x**3 for x in range(1,25,2)] # Cubes of odd numbers
  \end{lstlisting}  \pause
  \begin{itemize}
  \item Use an \lstinline{if} statement to check if 125 is in \lstinline{my_list}, and print a message indicating the result.
  \item Write a statement that checks and prints whether \lstinline{my_list[3]} is divisible by 3, and if not, print the remainder.
  \item Use a list comprehension to create a list of all numbers in \lstinline|my_list| that are not divisible by 5.
 \end{itemize}\pause
 \begin{onlyenv}<beamer>
  \begin{columns}[T]
    \column{0.33\textwidth}
    \begin{lstlisting}[]
>>> if 125 in my_list:
...     print('The number 125 was found in my_list!')
    \end{lstlisting}
    \column{0.33\textwidth}
    \begin{lstlisting}[]
>>> if my_list[3] % 3 == 0:
...     print(f'{my_list[3]} is divisible by 3')
... else:
...     print(f'The remainder of {my_list[3]} and 3 is {my_list[3]%3}')
    \end{lstlisting}
    \column{0.33\textwidth}
    \begin{lstlisting}
>>> k = [n for n in my_list if not n % 5 == 0 ]
    \end{lstlisting}
  \end{columns}
\end{onlyenv}
\end{frame}

%% FUNCTIONS 
\section{Functions}
\subsection{Defining functions}
\againframe<2>{contents_prog1}

\begin{frame}[fragile]
  \frametitle{Functions in Python (1)}
  Functions are defined using the \lstinline{def} keyword followed by the function name and a list of parameters in parentheses. The function body starts after the colon:
  \begin{lstlisting}[language=Python, numbers=none]
>>> def greet(name):
...     print(f"Hello, {name}!")
  \end{lstlisting}\pause
  Call the function with the necessary arguments:
  \begin{lstlisting}[language=Python, numbers=none]
>>> greet("Alice")
Hello, Alice!
  \end{lstlisting}\pause
  Functions can return values using the \lstinline{return} keyword:
  \begin{lstlisting}[language=Python, numbers=none]
>>> def add(a, b):
...     return a + b
  \end{lstlisting}\pause
  Capture the return value in a variable, e.g. \lstinline|result|:
  \begin{lstlisting}[language=Python, numbers=none]
>>> result = add(2, 3)
>>> print(result)
5
  \end{lstlisting}
\end{frame}

\begin{frame}[fragile]
  \frametitle{Functions in Python (2)}
  Default argument values can be specified, making the argument optional:
  \begin{lstlisting}[language=Python, numbers=none]
>>> def greet(name, greeting="Hello"):
...     print(f"{greeting}, {name}!")
  \end{lstlisting}\pause
  Call the function with or without the optional argument:
  \begin{lstlisting}[language=Python, numbers=none]
>>> greet("Bob")
Hello, Bob!
>>> greet("Bob", "Buzz off")
Buzz off, Bob!
  \end{lstlisting}\pause
  Python supports functions with a variable number of arguments:
  \begin{lstlisting}[language=Python, numbers=none]
>>> def my_function(*args):
...     print(args)
  \end{lstlisting}\pause
  Call the function with a varying number of arguments:
  \begin{lstlisting}[language=Python, numbers=none]
>>> my_function(1, 2, 3, "Hello")
(1, 2, 3, "Hello")
  \end{lstlisting}
\end{frame}

\begin{frame}[fragile]
  \frametitle{Functions in Python (3)}
  Functions can also return multiple values (also >2)
  \begin{lstlisting}[language=Python, numbers=none]
>>> def statistics(numbers):
...     return max(numbers), min(numbers)
  \end{lstlisting}\pause
  Let's call the function with some list:
  \begin{lstlisting}[language=Python, numbers=none]
>>> numlist = [94,12,6,19,33,14,81,56,43,22]
>>> print(statistics(numlist))
(94, 6)
  \end{lstlisting}\pause
  Store the elements in separate variables:
  \begin{lstlisting}[language=Python, numbers=none]
>>> a,b = statistics(numlist)
>>> print(f'{a=}, {b=}')
a=94, b=6
  \end{lstlisting}
\end{frame}

{\nologo
\begin{frame}[fragile]
  \frametitle{Functions in Python (4)}
  \begin{itemize}
    \item Functions are very useful for \emph{abstraction}:
    \begin{itemize}
      \item You can compartmentalize and 'hide' complex pieces of code\
      \item Retain flexibility through arguments
      \item You can reuse often used pieces of code, limiting copy-paste of code
    \end{itemize}
    \item Extending functionality or fixing bugs is done in 1 place\pause
    \item \bfseries{Documentation is crucial!}\pause
  \end{itemize}
  \begin{lstlisting}[language=Python, numbers=none]
>>> def statistics(numbers):
...     """Return the maximum and minimum of a list of numbers
...        Function arguments:
...        numbers: list of numbers"""
...     return max(numbers), min(numbers)
  \end{lstlisting}\pause
  \begin{lstlisting}[language=Python, numbers=none]
>>> help(statistics)
  \end{lstlisting}
  \begin{lstlisting}[style=PyOutput,keywords={statistics},keywordstyle=\bfseries]
Help on function statistics in module __main__:

statistics(numbers)
Return the maximum and minimum of a list of numbers
Function arguments:
numbers: list of numbers
  \end{lstlisting}
\end{frame}
}

{\nologo
\begin{frame}[fragile]
  \frametitle{Passing arguments by value or reference?}
  Recall that certain Python variables are assigned by value, and others reference; the same goes for passing arguments \footnote{\href{https://k0nze.dev/posts/python-copy-reference-none/}{https://k0nze.dev/posts/python-copy-reference-none/}}:
  \begin{itemize}
    \item Primitive or immutable data types are \emph{passed by value}; the value is copied and changes made inside the function will not affect the values stored in the variables passed to the function.
    \item Mutable data types (e.g. list, set, dict) are passed to a function \emph{by reference}; changes that are made inside the function do affect the values outside of the function.
  \end{itemize}
  \begin{columns}[t]
    
    \column{0.4\textwidth}
    Consider the function:

    \begin{lstlisting}
def func(x, y):
    x = x - 1 # Subtract 1
    y.pop()   # Remove last item
    \end{lstlisting}\pause
    \column{0.6\textwidth}
    % We call the function using an immutable data type (\lstinline|int|) and a mutable data type (\lstinline|list|):
    \begin{lstlisting}
i = 1
l = ['a', 'b']

func(i, l)

print(i, l)
    \end{lstlisting}
    \begin{lstlisting}[style=PyOutput]
1, ['a']      
    \end{lstlisting}
  \end{columns}
\end{frame}
}


\subsection{Scope}
\begin{frame}[fragile]
  \frametitle{Scope of functions and variables in Python}

  In Python, the scope of a variable refers to the regions of a program where that variable is accessible. Understanding the scope of variables helps to avoid bugs and maintain a clean codebase. The scopes in Python are categorized as follows:

  \begin{description}
    \item[Local Scope] Variables defined inside a function are in the local scope of that function. They can only be accessed within that function.
    \item[Enclosing Scope] In the case of nested functions, a function will have access to the variables of the functions it is nested within.
    \item[Global Scope] Variables defined at the top-level of a script are global and can be accessed by all functions in the script, unless overridden within a function.
    \item[Built-in Scope] Python has a number of built-in identifiers that should not be used as variable names as they have special significance. Examples include \textit{print}, \textit{list}, \textit{dict}, etc.
  \end{description}

  % Python follows the LEGB rule for name resolution: Local, Enclosing, Global, Built-in.
\end{frame}

{\nologo
\begin{frame}[fragile]
  \frametitle{Examples Variable Scope}
  \begin{columns}[T]
    \column[T]{0.5\textwidth}
      \textbf{1. Local Scope}:
    \column[T]{0.5\textwidth}
      \begin{lstlisting}[language=Python]
def my_func():
local_var = 100  # Local scope
    print(local_var)
      \end{lstlisting}
  \end{columns}
  {\color{scharlaken}\noindent\rule{\textwidth}{0.4pt}}
  \begin{columns}    
  \column[T]{0.5\textwidth}
    \textbf{2. Enclosing Scope}:
  \column[T]{0.5\textwidth}
    \begin{lstlisting}[language=Python]
def outer_func():
    outer_var = 200  # Enclosing scope
          
    def inner_func():
        print(outer_var)
    
    inner_func()
    \end{lstlisting}
  \end{columns}
  {\color{scharlaken}\noindent\rule{\textwidth}{0.4pt}}
  \begin{columns}[T]
    \column[T]{0.5\textwidth}
    \textbf{3. Global Scope}:
    \column[T]{0.5\textwidth}
      \begin{lstlisting}[language=Python]
global_var = 300  # Global scope

def another_func():
    print(global_var)
      \end{lstlisting}
    \end{columns}
    {\color{scharlaken}\noindent\rule{\textwidth}{0.4pt}}
    \begin{columns}[T]
      \column[T]{0.5\textwidth}
        \textbf{4. Built-in Scope}:
      \column[T]{0.5\textwidth}
        \begin{lstlisting}[language=Python]
print(max, min, len, str, int, list)
        \end{lstlisting}
    \end{columns}
\end{frame}
}

\begin{frame}[fragile]
  \frametitle{Exercise Variable Scope}
  Investigate the behavior of the following nested functions and variables with the same name:
  \begin{lstlisting}[language=Python]
def outer_func():
    outer_var = 200

    def inner_func():
        outer_var = 500
        print(outer_var)

    inner_func()
  \end{lstlisting}
\end{frame}

\begin{frame}[fragile]
  \frametitle{Lambda functions}
  If you do not want to create a new function block, you can create an \emph{lambda function}:
  Lambda functions are small, anonymous functions that can be instantiated in a single line, or even as an argument to a function.
  \begin{lstlisting}[language=Python]
from math import exp
f = lambda x: x**2 + exp(x)
  \end{lstlisting}
  \pause
  \begin{itemize}
    \item \lstinline[language=Python]|f|: the name of the function
    \item \lstinline[language=Python]|lambda|: used to define the inline function
    \item \lstinline[language=Python]|x|: the input argument (can be multiple, comma separated)
    \item \lstinline|:|: colon indicating the function definition will start
    \item \lstinline[language=Python]|x**2 + exp(x)|: the actual function
  \end{itemize}\pause
  \begin{lstlisting}[language=Python]
xsqr_exp = [f(x) for x in range(5)]
print(xsqr_exp)
  \end{lstlisting}
  \begin{lstlisting}[language=Python]
[1.0, 3.718281828459045, 11.38905609893065, 29.085536923187668, 70.59815003314424]
  \end{lstlisting}
\end{frame}




\subsection{The random module}
\begin{frame}[fragile]
  \frametitle{The random module (1)}
  Random number generators and sampling tools are available through the \href{https://docs.python.org/3/library/random.html}{random module}. A few examples for \href{https://docs.python.org/3/library/random.html#functions-for-integers}{integers and sequences}:
  \begin{lstlisting}
import random as rnd

# Random integers
random_integers = [rnd.randint(0,10) for i in range(10)]
print(f'{random_integers = }')

# Sample from given population
my_range = range(12) # [0, 1, 2, ..., 11]
select_from_range = rnd.sample(my_range,8)
print(f'{select_from_range = }') # Selected 8 elements

# Choose 1 element from list
days_of_week = ['Monday','Tuesday','Wednesday','Thursday','Friday','Saturday','Sunday']
day = rnd.choice(days_of_week)
print(f"I've chosen {day} as my lucky day!")
  \end{lstlisting}\pause

  \begin{lstlisting}[style=PyOutput]
random_integers = [6, 6, 2, 5, 5, 8, 3, 7, 3, 4]
select_from_range = [2, 3, 10, 1, 6, 4, 8, 5]
I've chosen Wednesday as my lucky day!
  \end{lstlisting}
\end{frame}

\begin{frame}[fragile]
  \frametitle{The random module (2)}
  Examples for \href{https://docs.python.org/3/library/random.html#real-valued-distributions}{real valued distributions}:
  \begin{lstlisting}
import random as rnd

# Random number (uniform distribution) 0 < x < 1
x = [rnd.random() for i in range(5)]
print(x)

# Random number (uniform distribution) between given bounds
x = [rnd.uniform(1,3) for i in range(5)]
print(x)

# Random number from a Gauss distribution (mu=0, sigma=2)
x = [rnd.gauss(0,2) for i in range(5)]
print(x)
  \end{lstlisting}\pause

  \begin{lstlisting}[style=PyOutput,basicstyle=\tiny\ttfamily]
[0.6697878114597362, 0.4136014290205997, 0.5108247513505662, 0.44260043089156076, 0.9902269207988261]
[2.4749381508841077, 2.943448233960596, 2.516639180020423, 1.0481550073898795, 1.961356325508141]
[1.8199856149229392, 2.000097897306016, -2.4604868187736026, -0.46836605162997846, -2.5069012642608803]
  \end{lstlisting}
\end{frame}
  
\begin{frame}[fragile,label=practice_random]
  \frametitle{Practice}
  \begin{itemize}
    \item Create a \emph{function} that returns a list of $N$ dice throws (cubic dice, values 1-6)
    \item Throw the dice many times
    \item Print for each value how often it has been thrown
  \end{itemize} \pause
  \begin{onlyenv}<beamer>
    \begin{lstlisting}
import random as rnd

def throw_dice(N):
    return [rnd.randint(1,6) for _ in range(N)]

throws = throw_dice(40)
print(throws)

for i in range(1,7):
    n = len([t for t in throws if t==i])
    print(f"Value {i} was thrown {n} times")
        \end{lstlisting}\pause
  \end{onlyenv}
  \begin{lstlisting}[style=PyOutput,basicstyle=\tiny\ttfamily]
[5, 1, 1, 4, 6, 3, 3, 1, 5, 6, 1, 1, 1, 1, 5, 5, 2, 1, 1, 2, 2, 5, 5, 4, 6, 1, 3, 5, 6, 3, 1, 5, 6, 2, 3, 1, 6, 3, 2, 1]
Value 1 was thrown 13 times
Value 2 was thrown 5 times
Value 3 was thrown 6 times
Value 4 was thrown 2 times
Value 5 was thrown 8 times
Value 6 was thrown 6 times
  \end{lstlisting}
\end{frame}
